% =========================================================
% VI. RESULTS AND ANALYSIS
% =========================================================
\section{Results and Analysis}
\label{sec:results_analysis}

\subsection{Overall Results and the Closed-Book Baseline}

All four runs produced 250/250 successful predictions with no pipeline errors
and shared the same model and generation settings. Table~\ref{tab:main_results}
reports end-to-end quality; retrieval/path cells are discussed only for systems
that expose those components.

\begin{table}[t]
\centering
\caption{Overall performance on 250 questions (percent except time).}
\label{tab:main_results}
\scriptsize
\setlength{\tabcolsep}{1.35pt}
\begin{tabular}{@{}lrrrrrrr@{}}
\hline
\textbf{Method}
& \textbf{Acc.}
& \textbf{B.Acc.}
& \textbf{M-F1}
& \textbf{Ans.F1}
& \textbf{R-L}
& \textbf{Cit.F1}
& \textbf{Sec.}\\
\hline
Closed-book LLM & 0.00 & 0.00 & 0.00 & 12.98 & 11.27 & 0.00 & 0.841\\
Vanilla RAG & 55.60 & 59.89 & 67.60 & 23.28 & 21.45 & 49.37 & 2.308\\
Causal Path RAG & 90.80 & 88.15 & 84.36 & \textbf{36.53} & \textbf{35.82} & 58.70 & 2.107\\
Full LegalSCM & \textbf{96.80} & \textbf{95.98} & \textbf{93.31} & 36.32 & 35.73 & \textbf{59.30} & 3.324\\
\hline
\end{tabular}
\end{table}

The 0.00\% closed-book decision score is not a failed run or an evaluator
parsing error. All 250 successful generations returned the valid label
\textsc{uncertain}. Since the gold set contains 230 \textsc{supported}, 20
\textsc{reject-direct-claim}, and no \textsc{uncertain} item, exact-label
accuracy is zero. The result therefore measures \emph{complete abstention under
this conservative shared prompt/model configuration}; it must not be read as
zero general legal competence. Conversely, an always-\textsc{supported}
classifier would obtain 92\% raw accuracy on this imbalanced set, which is why
Balanced Accuracy and Macro-F1 are reported alongside Accuracy.

\subsection{Paired Effect of Selective Verification}

Full LegalSCM answers 242 questions correctly versus 227 for Causal Path RAG.
Among paired outcomes, Full alone is correct on 17 samples, Causal Path RAG
alone on 2, both are correct on 225, and both are wrong on 6. The accuracy gain
is therefore 6.0 percentage points, with paired bootstrap 95\% CI
$[2.8,9.6]$ points, exact McNemar $p=0.000729$, and Holm-adjusted
$p=0.002914$. Macro-F1 increases descriptively by 8.94 points, from 84.36\% to
93.31\%.

The gain is decision-specific. Relative to Causal Path RAG, Full changes Answer
F1 by $-0.21$ points (95\% CI $[-1.83,1.43]$), ROUGE-L by $-0.09$ points
(95\% CI $[-1.71,1.56]$), and Citation F1 by $+0.59$ points (95\% CI
$[-0.72,2.03]$). All three Holm-adjusted $p$-values equal 1.0. Thus the data
support improved structured decisions, not a general improvement in answer
wording or citations.

\subsection{Route, Commit, and Error Audit}

Table~\ref{tab:route_audit} reports independent runtime counters. The query
router sends 230 samples to factual LegalSCM validation and 20 direct claims to
grounded edge checking. No query requests the mediator-intervention route.

\begin{table}[t]
\centering
\caption{Full LegalSCM route and commit audit.}
\label{tab:route_audit}
\small
\begin{tabular}{lr}
\hline
\textbf{Audit quantity} & \textbf{Count}\\
\hline
Factual LegalSCM route & 230\\
Grounded direct-edge route & 20\\
Verifier committed & 230\\
No-primary-path safe fallback & 20\\
Executed $do(X_M=0)$ intervention & 0\\
\hline
\end{tabular}
\end{table}

All 17 Full-only gains occur when the verifier commits: 14 arise from factual
rule validation and 3 from direct-edge topology. In the 20 safe-fallback cases,
Full and Causal Path RAG make the same decision on 19 samples; 13 are jointly
correct and 6 jointly wrong. The remaining case is correct only for Causal Path
RAG. Hence the fallback contributes no Full-only gain.

Full makes eight decision errors. Seven occur after no-primary-path fallback,
primarily because first-hop retrieval does not expose a usable primary path.
The remaining error is a direct-claim endpoint-grounding error. In three error
cases the generated prose also contradicts the structured label, identifying a
separate generation-consistency issue.

The direct-claim counterexample subset yields 95\% accuracy for Full and 85\%
for Causal Path RAG. However, all 20 samples are evaluated by direct-edge
topology and none executes $do(X_M=0)$. This subset supports the value of
query-aware direct-claim checking, not the effectiveness of structural mediator
intervention.

\subsection{Retrieval, Path Ranking, and Runtime}

Causal Path RAG and Full share the same retrieval output. Article Recall@5,
Rule Recall@5, and Event Recall@5 are 75.40\%, 71.20\%, and 74.79\%.
Top-1 Exact Path is 38.56\%, while Oracle Exact Path is 75.42\%. The 36.86-point
gap shows that the gold path is often in the candidate set but ranked below the
first position; retrieval/path equality also localizes the Full--Causal
decision difference to the selective verification layer rather than retrieval.

Full requires 3.324 seconds per successful question versus 2.107 seconds for
Causal Path RAG, a $1.58\times$ average-time increase. The measured trade-off is
therefore a statistically reliable 6.0-point decision gain at approximately
58\% additional per-sample runtime.
